%%% Kapitola 4: Testování

\chapter{Testování}
\label{chap:testovani}

Cílem testování bylo ověřit, že navržená aplikace splňuje funkční požadavky
formulované v~kapitole~\ref{chap:specifikace} -- zejména že dokáže přihlásit
uživatele přes Trezor, správně sestavit a~odeslat singlesig transakci a~že celý
multisig workflow od importu deskriptoru přes postupné podepisování až po
broadcast funguje proti reálnému Bitcoinovému uzlu. V~této kapitole popisujeme
zvolený přístup k~testování, jeho odůvodnění, použité prostředí, provedené
scénáře a~nakonec chyby, které jsme díky testování odhalili.

%%% -----------------------------------------------------------------------
\section{Zvolený přístup}
\label{sec:test-pristup}

Při návrhu testovací strategie jsme zvažovali tři varianty.

První možností byla \textbf{automatizovaná sada unit testů} pokrývající
jednotlivé komponenty backendu (parsing deskriptorů, derivaci adres,
sestavování PSBT). Tento přístup poskytuje rychlou zpětnou vazbu při
změnách kódu a~je standardem pro produkční software. V~kontextu naší
aplikace by však unit testy izolovaly právě ty části, kde jsme reálně
nehledali chyby -- kryptografická logika je delegována na knihovnu
bitcoinj~\cite{BitcoinJ}, derivace probíhá podle dobře definovaných BIP
standardů a~lze ji ověřit proti existujícím peněženkám. Naopak by neodhalily
problémy v~integraci s~Trezorem, protože ten je v~unit testech nahrazen
mockem, který se na rozdíl od skutečného firmware chová tolerantně.

Druhou možností byly \textbf{end-to-end testy s~emulátorem Trezoru}, který
SatoshiLabs poskytuje pro vývojové účely. Tento přístup by umožnil
automatizovat i~podepisovací krok. Komunikace s~emulátorem však vyžaduje
implementaci části protokolu Trezor Connect, což je přesně to, čemu se
naše architektura záměrně vyhýbá delegováním na Trezor Suite Mobile (viz
sekci~\ref{sec:analyza-trezor}). Pro testovací sadu by tak bylo nutné
vyvinout paralelní komunikační vrstvu, která by se v~samotné aplikaci
vůbec nepoužila. To jsme v~rámci bakalářské práce vyhodnotili jako úsilí
neúměrné přínosu.

Zvolili jsme proto třetí variantu: \textbf{manuální integrační testování
proti Bitcoin testnetu} s~reálnou hardwarovou peněženkou Trezor.
Tento přístup ověřuje právě ty integrační body, kde reálně vznikají chyby --
komunikaci mezi mobilní aplikací, Trezor Suite Mobile a~hardwarovým
zařízením, korektnost sestavených PSBT z~pohledu firmware Trezoru
a~validitu výsledných transakcí z~pohledu Bitcoinové sítě. Cenou je nutnost
manuálně procházet všechny scénáře při každé větší změně, což akceptujeme
vzhledem k~omezenému rozsahu funkcionality. Doplnění automatizované testovací
sady považujeme za vhodný směr dalšího vývoje.

%%% -----------------------------------------------------------------------
\section{Testovací prostředí}
\label{sec:test-prostredi}

Veškeré testování probíhalo proti síti testnet4. Testnetové prostředky
(tBTC) jsme získávali z~veřejně dostupných faucetů -- webových služeb, které
bezplatně distribuují malé množství testnetových mincí vývojářům pro účely
testování.

Backend byl spuštěn lokálně pomocí Docker Compose~\cite{Docker} -- jednou
spuštěný stack obsahoval všech šest mikroslužeb popsaných
v~sekci~\ref{sec:architektura},
API Gateway, tři instance PostgreSQL~\cite{PostgreSQL} a~síťové připojení na
veřejné API Mempool.space~\cite{MempoolAPI} pro testnet4. Mobilní aplikace
běžela na fyzickém Android zařízení se soubežně nainstalovanou aplikací
Trezor Suite Mobile, která zprostředkovává komunikaci s~hardwarovou peněženkou
podle deeplink protokolu popsaného v~sekci~\ref{subsec:trezor-connect}.
K~testování jsme využili jednu hardwarovou peněženku Trezor s~inicializovaným
seedem.

Pro multisig scénáře 2-of-3 vyvstává otázka, jak simulovat tři nezávislé
cosignery, když máme k~dispozici pouze jedno fyzické zařízení. Zvolili jsme
přístup, kdy na témže Trezoru používáme tři různé účty podle BIP-48~\cite{BIP48}
s~derivačními cestami \texttt{m/48'/1'/0'/2'}, \texttt{m/48'/1'/1'/2'}
a~\texttt{m/48'/1'/2'/2'}. Z~pohledu kryptografického protokolu jsou tyto účty
plnohodnotnými nezávislými klíčovými páry -- BIP-32 garantuje, že odvozené
klíče různých větví spolu nesouvisí způsobem, který by šel z~veřejných dat
detekovat. Aplikace přitom mezi účty přepíná na úrovni session a~Trezor pro
každý podpis vyžaduje samostatné fyzické potvrzení, takže z~hlediska workflow
se postup neliší od reálného nasazení s~více zařízeními. Jediné, co tento
přístup neověří, je souběžné podepisování více uživateli v~čase; v~našem
sériovém workflow však k~souběhu beztak nedochází, protože PSBT se mezi
cosignery předává postupně.

Multisig deskriptor jsme vytvářeli v~aplikaci Sparrow Wallet~\cite{SparrowWallet},
do které jsme z~Trezoru exportovali tři xpub klíče, nakonfigurovali peněženku
2-of-3 a~deskriptor ve formátu \texttt{wsh(sortedmulti(2,...))} jsme zkopírovali
do importního dialogu naší aplikace. Sparrow zároveň sloužil jako referenční
peněženka -- adresy odvozené naší aplikací jsme porovnávali s~adresami, které
pro stejný deskriptor zobrazil Sparrow, a~v~prohlížeči blockchainu Mempool.space
jsme nezávisle ověřovali stav transakcí.

%%% -----------------------------------------------------------------------
\section{Provedené scénáře}
\label{sec:test-scenare}

Testovací scénáře jsme odvodili z~případů užití definovaných
v~sekci~\ref{sec:use-cases}. Pro každý use case existuje odpovídající
manuální scénář; níže je popisujeme stručně, pouze z~hlediska toho, co
ověřují a~co je u~nich kritické.

\textbf{Singlesig workflow.} Začínáme přihlášením uživatele přes Trezor
(UC1). Aplikace vyžádá z~hardwarové peněženky deset xpub klíčů pro účty
\texttt{m/84'/1'/0'} až \texttt{m/84'/1'/9'}, backend pro každý xpub provede
account discovery dle BIP-44~\cite{BIP44} a~vrátí seznam aktivních
singlesig peněženek. Tento scénář ověřuje, že deeplink komunikace s~Trezor
Suite Mobile funguje a~že JWT token vydaný službou \texttt{auth-service}
umožní následný přístup k~ostatním službám.

Dalším scénářem je odeslání singlesig transakce typu P2WPKH (UC2). Zadáme
testnetovou adresu, částku a~úroveň poplatku, aplikace nás přes deeplink
přesměruje do Trezor Suite Mobile, kde transakci podepíšeme na zařízení.
Po návratu callback URL aplikace odešle serializovanou transakci na
backend, který provede broadcast přes \texttt{blockchain-service}. Úspěšnost
ověřujeme jednak hlášením z~aplikace (zobrazení txid), jednak vyhledáním
této transakce v~Mempool.space. Variantou tohoto scénáře je test funkce
coin control (UC5), kdy před potvrzením ručně vybereme konkrétní UTXO --
ověřujeme, že výsledná transakce na blockchainu obsahuje právě jeden vstup
odpovídající zvolenému UTXO a~že odhad poplatku se přepočítal podle nového
počtu vstupů.

Příjem prostředků (UC z~user story~US4) testujeme tak, že v~aplikaci
zobrazíme přijímací adresu, z~externího faucetu na ni odešleme malou částku
tBTC a~ověřujeme, že po potvrzení transakce se zůstatek v~aplikaci aktualizuje
a~v~historii se objeví příchozí transakce se správnou klasifikací směru
(\textit{received}, nikoli \textit{sent}).

\textbf{Multisig workflow.} Po importu deskriptoru (UC3) ověřujeme, že
aplikace zobrazí novou peněženku s~informací ,,2-of-3``, že odvozené adresy
začínají prefixem \texttt{tb1q} a~jsou shodné s~adresami, které pro tentýž
deskriptor zobrazí Sparrow Wallet, a~že zobrazený zůstatek odpovídá
prostředkům, které jsme na multisig adresy předem odeslali z~faucetu.

Hlavním scénářem multisig části je vytvoření, podepsání a~broadcast 2-of-3
transakce (UC4 + UC6). Postup probíhá ve čtyřech krocích: nejprve jako
cosigner~1 (\texttt{m/48'/1'/0'/2'}) zadáme cílovou adresu a~částku, backend
sestaví PSBT pro multisig vstupy a~aplikace nás přes deeplink odešle k~podpisu
na Trezor. Trezor v~multisig režimu vrací pouze pole podpisů, nikoli celou
serializovanou transakci -- tyto podpisy backend přidá do PSBT na pozici
odpovídající aktuálnímu cosignerovi. Stav PSBT přejde z~\texttt{pending}
(0~podpisů) na 1~ze~2 a~aplikace zobrazí informaci, že transakce čeká na
dalšího cosignera. Ve druhém kroku v~nastavení přepneme aktivní účet na
cosignera~2 (\texttt{m/48'/1'/1'/2'}), v~seznamu rozpracovaných PSBT vybereme
připravovanou transakci a~znovu projdeme deeplink podpisem. Backend po
přijetí druhého podpisu uloží podpisy na~pozice odpovídající
BIP-67~\cite{BIP67} řazení a~vrátí aktualizovaný stav PSBT; aplikace
detail obnoví a~místo tlačítka \emph{Sign with Trezor} zobrazí
\emph{Broadcast Transaction}. Po~jeho stisknutí aplikace zavolá
broadcast endpoint, který transakci propaguje přes
\texttt{blockchain-service} do~sítě. Úspěch ověřujeme v~Mempool.space,
kde u~vstupů vidíme očekávaný witness
obsahující dva podpisy a~kompletní redeem skript ve tvaru
\texttt{OP\_2~<pk1>~<pk2>~<pk3>~OP\_3 OP\_CHECKMULTISIG}.

V~rámci tohoto scénáře ověřujeme i~speciální případy, které byly v~průběhu
vývoje zdrojem chyb -- zejména transakci s~více vstupy z~různých multisig
adres (každý vstup může mít jiné lexikografické pořadí cosignerů, takže
podpisy nelze řadit globálně, viz~sekci~\ref{sec:test-chyby}) a~transakci
s~change výstupem, u~níž Trezor musí na svém displeji označit change jako
vlastní výstup, nikoli jako platbu třetí straně.

%%% -----------------------------------------------------------------------
\section{Nalezené chyby}
\label{sec:test-chyby}

Manuální testování splnilo svůj hlavní účel -- odhalilo několik integračních
chyb, které by se v~unit testech s~mockovaným Trezorem nemohly projevit.
Dvě z~nich považujeme za dostatečně netriviální, abychom je popsali
podrobněji.

\textbf{Chybné řazení podpisů ve witness (BIP-67).}
Při prvním pokusu o~broadcast 2-of-3 transakce s~více vstupy Bitcoinová síť
transakci odmítla jako nevalidní, ačkoli oba cosigneři transakci úspěšně
podepsali. Vyšetřování ukázalo, že podpisy byly ve witness datech umístěny
ve špatném pořadí. Příčinou byla skutečnost, že naše původní implementace
řadila pozice cosignerů jednou globálně podle pořadí xpub klíčů v~deskriptoru,
zatímco BIP-67~\cite{BIP67} vyžaduje řazení na úrovni odvozených (child)
veřejných klíčů, a~to pro každou adresu samostatně. Protože child klíče
jednotlivých cosignerů odvozené pro různé adresy mohou mít různé vzájemné
lexikografické pořadí, globální strategie fungovala jen tehdy, když se
shodou okolností shodovala s~lokálním pořadím u~první adresy. Oprava
spočívala v~tom, že komponenta sestavující witness pro každý vstup nezávisle
odvodí child klíče všech cosignerů, seřadí je dle BIP-67 a~podpis daného
cosignera umístí na pozici odpovídající v~tomto lokálním seřazení. Od této
opravy procházejí všechny multisig scénáře včetně transakcí s~více vstupy.

\textbf{Chybějící \texttt{PSBT\_IN\_NON\_WITNESS\_UTXO}.}
Po upgradu firmware Trezoru začaly i~singlesig transakce, které dosud
fungovaly bez problémů, končit chybou ,,Transaction has changed during
signing``. Dohledali jsme, že firmware Trezoru od verze 2.4 vyžaduje, aby
PSBT pro každý vstup obsahoval kompletní předchozí transakci v~poli
\texttt{PSBT\_IN\_NON\_WITNESS\_UTXO} -- a~to i~u~nativních SegWit vstupů
typu P2WPKH, kde to BIP-174 striktně nevyžaduje. Důvodem je, že firmware
chce nezávisle ověřit částky vstupů a~předejít útoku, při kterém by se mu
podsunula falešná hodnota poplatku~\cite{TrezorFirmwareDocs}. Naše původní
implementace toto pole vynechávala, protože pro P2WPKH stačilo
\texttt{PSBT\_IN\_WITNESS\_UTXO} obsahující pouze utrácené UTXO.
Oprava vyžadovala dvě změny: ve službě \texttt{blockchain-service} jsme
přidali endpoint \texttt{/tx/\{txid\}/hex} vracející serializovanou
předchozí transakci a~ve službě \texttt{psbt-service} při sestavování PSBT
pro každý vstup tento raw tx fetchneme a~zapíšeme do odpovídajícího pole.

Kromě těchto dvou jsme během testování odhalili a~opravili i~další chyby.
Významnou z~nich bylo chybějící pole
\texttt{PSBT\_OUT\_BIP32\_DERIVATION} u~change výstupu. Toto pole podle
BIP-174~\cite{BIP174} sděluje podepisujícímu zařízení derivační cestu
výstupu, díky čemuž může ověřit, že change adresa patří téže peněžence
a~zobrazit ji uživateli jako vnitřní převod, nikoliv jako platbu třetí
straně. Bez tohoto pole Trezor při podepisování zobrazoval change výstup
jako samostatnou platbu na cizí adresu, což uživatele mátlo -- na displeji
viděl dvě odchozí platby místo jedné platby a~jednoho vrácení změny.
Oprava spočívala v~doplnění derivační cesty change adresy do příslušného
pole při sestavování PSBT.

Zajímavým příkladem toho, jak nová funkce odhalí skrytou chybu, bylo
opravení derivace přijímací adresy. Služba \texttt{explorer-\allowbreak service}
v~metodě \texttt{getNext\-Receive\-Address} vracela poslední \emph{použitou}
adresu místo první \emph{nepoužité} -- všechny příchozí platby tak směřovaly
na tutéž adresu, což porušuje zásadu neopakování adres (každá transakce
by měla používat novou adresu kvůli ochraně soukromí). Chybu jsme odhalili
díky funkci rekurzivního procházení vstupů v~detailu transakce
(viz sekci~\ref{sec:user-detail-tx}): při proklikávání předchozích transakcí
jsme si všimli, že více nezávislých UTXO odkazuje na tutéž přijímací adresu.
Oprava spočívala v~přepsání logiky tak, aby služba prohledala adresy od
nejnižšího indexu a~vrátila první, která dosud nemá žádnou aktivitu
v~blockchainu; pokud jsou všechny předem odvozené adresy použité, služba
automaticky derivuje novou adresu na dalším indexu.

Dále jsme opravili počítadlo virtuální velikosti, které pro multisig vstupy
používalo hodnotu 68~vB platnou pro singlesig P2WPKH a~vedlo k~podstatnému
podhodnocení odhadovaného poplatku u~multisig transakcí.

%%% -----------------------------------------------------------------------
\section{Známá omezení}
\label{sec:test-omezeni}

Zvolený přístup k~testování má několik omezení, která je vhodné explicitně
uvést. Prvním je již zmíněná simulace tří cosignerů na jednom fyzickém
zařízení -- ačkoli je z~kryptografického pohledu plnohodnotná, neumožňuje
zachytit potenciální problémy specifické pro produkční nasazení s~více
zařízeními (například chyby v~enkódování PSBT, které by se projevily pouze
mezi různými verzemi firmware). Druhým omezením je zaměření testování výhradně na testnet -- ačkoli
aplikace podporuje přepínání mezi testnetem a~mainnetem přímo
z~přihlašovací obrazovky, veškeré testovací scénáře jsme prováděli pouze
na testnetu. Provoz na mainnetu jsme neověřovali.
Konečně, neměřili jsme chování aplikace pod zátěží
ani její odolnost vůči síťovým výpadkům -- pro mobilní peněženku s~jednotkami
uživatelů jsou tyto charakteristiky druhořadé, ale v~širším nasazení by
si zasloužily samostatné ověření.
